\documentclass{situationist-preprint}

\title{Hope and Probability}
\msid{EXI-15-b}
\msfamily{EXI-15}
\msdiscipline{Meaning, Death and the Absurd}
\mstype{Companion manuscript}
\mskeywords{hope; probability; credence; wishful thinking; self-fulfilling attitudes; virtue}

\begin{document}
\maketitle

\begin{abstract}
EXI-15-a argues that hope is a virtue when it is proportioned to probability without collapsing into it: the strength of the hopeful stance must increase with credence, other things equal, but must depend also on the value of the outcome and on the agent's power to bring it about. We develop the relation between hope and probability in three directions. We describe how the stance should behave at the extremes of credence, where hope passes into expectation or into mere wishing. We distinguish the legitimate way in which hope can raise the probability of what is hoped for, by changing what the agent does, from the illegitimate way in which it can raise credence directly. And we answer the objection that the virtue so described is nothing but expected-value reasoning.
\end{abstract}

\section{The relation}

On the account of EXI-15-a, an agent hopes that $p$ if she desires $p$, has a credence in $p$ strictly between zero and one, and takes up a stance that treats the possibility of $p$ as a reason for planning, attention and effort (Definition~2.1 of EXI-15-a; \citealp{Bovens1999,Martin2014}). The credence is fixed by evidence (Principle~P1 of EXI-15-a); the stance is under the agent's control (Principle~P2). The virtue consists in proportion without collapse (Definition~3.1 of EXI-15-a): writing $c$ for the credence, $v$ for the value of $p$ to the agent and $e$ for her capacity to make $p$ more likely, the strength $s$ of the stance increases with $c$ when $v$ and $e$ are held fixed, and is not determined by $c$ alone.

The definition is qualitative. It does not specify a function $s(c,v,e)$, and we do not supply one. What it fixes is a pattern of comparative judgments: of two hopes alike in $v$ and $e$, the one with the higher credence warrants the stronger stance; of two hopes alike in $c$, the one with greater $v$ or $e$ may warrant the stronger stance. Lemma~3.2 of EXI-15-a condemns violations of the first comparison as a vice of excess; Lemma~3.3 condemns the refusal to make the second as deficiency.

\section{The extremes}

\begin{proposition}[Near certainty]\label{p:high}
As credence in $p$ approaches one, the hopeful stance passes into expectation: the agent plans on $p$ rather than for its possibility, and the specifically hopeful part of the stance has nothing further to do.
\end{proposition}

\begin{proof}
The stance treats the \emph{possibility} of $p$ as a reason to plan for it and to attend to the ways it might come about (Definition~2.1 of EXI-15-a). When $c$ is close to one, the rational agent plans on $p$ as she plans on what she expects; attention to the ways $p$ might come about adds nothing to the plan, since the plan already assumes it. Clause (a) of Definition~3.1 of EXI-15-a is satisfied, since the commitment of resources to $p$ is high, but the resources are now committed by expectation.
\end{proof}

\begin{proposition}[Remote possibility]\label{p:low}
As credence in $p$ approaches zero, a virtuous agent's stance towards $p$ approaches the minimal stance of wishing, unless the value of $p$ and her capacity to bring it about are great.
\end{proposition}

\begin{proof}
By clause (a) of Definition~3.1 of EXI-15-a, with $v$ and $e$ fixed, the stance weakens as $c$ falls. By clause (b), large $v$ and $e$ can sustain a substantial stance at low $c$, which is the case of hope against hope discussed in EXI-15-a. Where $v$ or $e$ is small, as in the hope of a lottery player who can do nothing to affect the draw, nothing sustains it, and the stance that remains is wishing.
\end{proof}

Neither proposition makes hope at the extremes a vice. The hope of a lottery player who buys a ticket and thinks no more of it is proportioned: weak, because $c$ and $e$ are small. The vice of excess is the player who reorganizes her life around the draw (Lemma~3.2 of EXI-15-a).

\section{Two ways hope can raise probability}

Hope sometimes helps to bring about what it hopes for. \citet{James1897} argued that in some cases faith in a fact can help to create the fact, and the observation is familiar in medicine and in collective action. There are two ways in which this might work, and the account of EXI-15-a sanctions only one.

\begin{proposition}\label{p:feedback}
When the agent's effort is among the causes of $p$, adopting the hopeful stance can raise the probability of $p$, and her credence may rationally rise in response. Adopting the stance cannot rationally raise her credence in any other way.
\end{proposition}

\begin{proof}
If $e$ is positive, the effort that the stance commits makes $p$ more likely; the agent's evidence then includes the fact that she is committed to the effort, and the credence that evidence supports is the probability of $p$ given her commitment. By Principle~P1 of EXI-15-a, her credence is fixed by her evidence, so it may rise. Any other route from stance to credence would make credence respond to desire rather than to evidence, which Principle~P1 excludes.
\end{proof}

The first route is how the hope of a patient who undertakes a demanding rehabilitation, or of a group that organizes for a cause, can be both rational and self-supporting. The second is wishful thinking. The distinction connects EXI-15 with EXI-18: where the function of an illusion of control or of unrealistic optimism can be served by hope, which by Proposition~\ref{p:feedback} raises probability only by changing what the agent does, hope is to be preferred, since it does not misrepresent (Remark~2.5 of EXI-18-a).

\section{Objections}

\begin{objection}[Expected value]
The virtue described is just expected-value reasoning: one invests in an outcome in proportion to its probability times its value, adjusted for what one's investment changes. Hope adds nothing.
\end{objection}

\begin{reply}
The comparative structure of Definition~3.1 of EXI-15-a is indeed close to that of expected-value reasoning, and that is a point in its favor: it shows that the virtue is a form of practical rationality and not a license to ignore the odds. But the hopeful stance includes more than choices among options. It includes where attention goes, which possibilities are rehearsed and prepared for, and the persistence of effort when evidence of progress is lacking \citep{McGeer2004}. Expected-value reasoning, applied to a given set of options, does not determine these, and it does not exclude the despairing agent who makes the right choices but ceases to attend to the possibility of success, and so fails to notice the openings that would raise it. The virtue governs that attention, and its standard is the one Definition~3.1 of EXI-15-a states.
\end{reply}

\begin{objection}[Hope corrupts credence]
In practice, hope does not stay on the first route of Proposition~\ref{p:feedback}. Those who hope strongly come to believe what they hope for, so that the virtue is unstable.
\end{objection}

\begin{reply}
That hope can corrupt credence shows that it can be exercised badly, as courage can pass into recklessness. A disposition whose stance is fixed by credence, and whose credence is fixed by evidence, is the disposition that resists this corruption; it is the virtue of Theorem~3.4 of EXI-15-a, and the corruption is one of the ways of falling short of it.
\end{reply}

\section{Conclusion}

The relation of hope to probability is one-directional at the level of attitude: credence constrains the stance and the stance does not constrain credence. It may be two-directional at the level of the world, since the stance can change what happens and the credence must then follow. Hope that keeps to this structure is proportioned without collapsing, and it is a virtue for the reasons EXI-15-a gives.

\begin{thebibliography}{9}\small
\bibitem[Bovens(1999)]{Bovens1999} Bovens, L. (1999). The value of hope. \emph{Philosophy and Phenomenological Research}, 59(3), 667--681.
\bibitem[James(1897)]{James1897} James, W. (1897). \emph{The Will to Believe and Other Essays in Popular Philosophy}. Longmans, Green, and Co.
\bibitem[Martin(2014)]{Martin2014} Martin, A.~M. (2014). \emph{How We Hope: A Moral Psychology}. Princeton University Press.
\bibitem[McGeer(2004)]{McGeer2004} McGeer, V. (2004). The art of good hope. \emph{The Annals of the American Academy of Political and Social Science}, 592, 100--127.
\end{thebibliography}

\end{document}
